\section{Phase 1 Interview Protocols}

Phase 1 used separate semi-structured interview guides for older adults and family caregivers. Interviews were combined with situated observation and medication-routine walkthroughs in participants' homes. Interviewers used the questions below flexibly and followed relevant responses with probes about recent or specific episodes. Where appropriate, participants were asked to demonstrate medication practices, artifacts, and technology use rather than relying only on general descriptions.

\subsection{Appendix A1: Older-Adult Interview Protocol}

\subsubsection{Phase 1: Environment and Artifact Observation}

\textbf{Observational focus:} Before beginning the interview, document the immediate medication environment, including medicine-storage locations, lighting and noise conditions, physical accessibility, and visible aids such as prescriptions, written notes, calendars, pillboxes, or other organizational strategies.

\subsubsection{Phase 2: Background and Medication Routine}

\begin{enumerate}
  \item \textbf{Can you tell me a little about yourself and how medicines are part of your everyday routine?}

  \textit{Probe: How many medicines do you usually take, and at what times of day?}
  \item \textbf{Can you walk me through how you normally take your medicines during a typical day?}

  \textit{Probe: What happens from the time you realize a medicine is due until you take it?}
  \item \textbf{How do you know which medicine to take and when?}

  \textit{Probe: Do you remember it yourself, check a prescription or written note, look at the medicine strip, or ask someone?}
  \item \textbf{Who usually takes responsibility for different parts of your medication routine?}

  \textit{Probe: Which parts do you normally manage yourself, and which parts does someone else help with?}
  \item \textbf{Has this arrangement always been the same?}

  \textit{Probe: How did family members become involved in the parts they now handle?}
\end{enumerate}

\subsubsection{Phase 3: Situated Medication Walkthrough}

\textbf{Observational focus:} Ask the participant to demonstrate a usual medication routine. Observe how medicines are located, identified, handled, and organized and what artifacts or workarounds support the process.

\begin{enumerate}[resume]
  \item \textbf{Could you show me how you would prepare and take one of your usual doses?}

  \textit{Probe: How do you identify the correct medicine and dose?}
  \item \textbf{Where do you usually keep your medicines? Why do you keep them there?}

  \textit{Probe: Does anyone else organize, move, or prepare them for you?}
  \item \textbf{Do you use anything to help you remember or organize your medicines?}

  \textit{Probe: Written notes, phone alarms, pillboxes, marked medicine strips, particular storage locations, or something else?}
  \item \textbf{Have you tried any method that you later stopped using?}

  \textit{Probe: What did you try, and why did you stop?}
  \item \textbf{What parts of your current routine work well, and what parts are difficult?}

  \textit{Probe: If you could change one part of the routine, what would you change?}
\end{enumerate}

\subsubsection{Phase 4: Disruptions, Missed Doses, and Handover}

\textbf{Observational focus:} Ask about specific recent episodes rather than only general difficulties. Attend to changes in routine, location, family availability, and other circumstances surrounding the episode.

\begin{enumerate}[resume]
  \item \textbf{Have you ever missed or delayed a dose?}

  \textit{Probe: Tell me about the most recent time. What was happening that day?}
  \item \textbf{How did you realize that the dose had been missed or delayed?}

  \textit{Probe: Did you notice yourself, or did someone else notice or remind you? What happened afterward?}
  \item \textbf{Are there situations when taking your medicines becomes more difficult than usual?}

  \textit{Probe: What happens when you are away from home, have visitors, attend a social occasion, sleep at a different time, or your usual schedule changes?}
  \item \textbf{What happens when the person who usually helps with your medicines is unavailable?}

  \textit{Probe: Tell me about the last time this happened. Did you manage differently yourself, or did someone else become involved?}
  \item \textbf{Have you ever taken the wrong medicine or taken a medicine twice by mistake?}

  \textit{Probe: How was it noticed, and what happened afterward?}
\end{enumerate}

\subsubsection{Phase 5: Prescription, Medicine Information, and Accessibility}

\textbf{Observational focus:} Where appropriate, invite participants to handle their own prescription, medicine strip, bottle, or other medication information while discussing how they use it.

\begin{enumerate}[resume]
  \item \textbf{How do you usually understand what is written on your prescription or medicine packaging?}

  \textit{Probe: Which parts can you read or recognize yourself?}
  \item \textbf{What do you do when something about a medicine or prescription is unclear?}

  \textit{Probe: Whom do you usually ask: a family member, doctor, pharmacist, or someone else?}
  \item \textbf{When a new medicine is prescribed, how do you learn when and how to take it?}

  \textit{Probe: Who usually explains it to you?}
  \item \textbf{Is anything about reading medicine labels, prescriptions, or instructions difficult for you?}

  \textit{Probe: Small text, medicine names, dosage instructions, language, or identifying similar-looking medicines?}
\end{enumerate}

\subsubsection{Phase 6: Technology and Device Use}

\textbf{Observational focus:} Where participants use a digital device, invite them to demonstrate relevant interactions. Observe shared or intermediated use rather than assuming that device ownership means independent use.

\begin{enumerate}[resume]
  \item \textbf{What kind of phone do you usually use?}

  \textit{Probe: Is it your own phone or a phone shared with other household members?}
  \item \textbf{What do you normally use your phone for?}

  \textit{Probe: Who helped you learn to use it or set it up?}
  \item \textbf{Do you currently use a phone, alarm, application, or other digital tool to help with medicines?}

  \textit{Probe: Could you show me how you use it?}
  \item \textbf{What feels easy or difficult when you use a phone or other technology?}

  \textit{Probe: Reading the screen, hearing alerts, typing, navigating menus, remembering steps, or something else?}
  \item \textbf{Does anyone in your family sometimes operate or configure a phone or application for you?}

  \textit{Probe: What do they usually help with?}
  \item \textbf{What would make a medication reminder or tracking tool easier for you to use?}

  \textit{Probe: Larger text, spoken information, simpler controls, or another form of support?}
\end{enumerate}

\subsubsection{Phase 7: Family Involvement and Medication Responsibility}

\begin{enumerate}[resume]
  \item \textbf{Who in your family is involved in your medicines?}

  \textit{Probe: What does each person usually do?}
  \item \textbf{How do family members know whether you have taken your medicine?}

  \textit{Probe: Do they ask you, watch the routine, look at the medicines, call you, or use some other way?}
  \item \textbf{What usually happens when someone reminds you about a medicine?}

  \textit{Probe: Can you tell me about a recent example?}
  \item \textbf{Are there parts of your medication routine that you prefer to manage yourself?}

  \textit{Probe: Which parts, and what makes those important for you to manage?}
  \item \textbf{When something about your medicines needs to change, how is that usually decided?}

  \textit{Probe: Who is involved in the discussion, and who normally acts on the change?}
  \item \textbf{Have you and a family member ever had different views about something related to your medicines?}

  \textit{Probe: Can you tell me about a specific occasion and what happened?}
  \item \textbf{How do you communicate with doctors or pharmacies about your medicines?}

  \textit{Probe: Do you usually communicate with them yourself, together with someone, or through another family member?}
\end{enumerate}

\subsubsection{Phase 8: Closing and Design Reflection}

\begin{enumerate}[resume]
  \item \textbf{If something could help with one part of your medication routine, which part would you most want help with?}

  \textit{Probe: Which parts would you prefer to continue managing yourself?}
  \item \textbf{If a reminder system could not tell whether you had taken a medicine, what do you think it should do next?}

  \textit{Probe: Should it remind you again, wait, or involve someone else?}
  \item \textbf{If someone else needed to be informed about a missed medicine, who should that person be?}

  \textit{Probe: How should that decision be made?}
  \item \textbf{Who should be able to change the settings of a medication-support system?}

  \textit{Probe: Should you know when someone else makes a change?}
  \item \textbf{Is there anything a medication-support system should never do?}
  \item \textbf{Is there anything about managing your medicines that we have not talked about but that you think is important for us to understand?}
\subsection{Appendix A2: Family-Caregiver Interview Protocol}
\end{enumerate}


The caregiver was interviewed as a participant in their own right. Questions focused on the caregiver's own work, judgments, and position within household medication routines. When relevant, interviewers distinguished between the caregiver's own experience, their account of the older adult's experience, and events they described as occurring jointly.

\subsubsection{Phase 1: Environment and Artifact Observation}

\textbf{Observational focus:} Document where medicines are stored, whether these spaces are managed primarily by the caregiver, the older adult, or jointly, and what medication aids are currently used or visibly abandoned. Note relevant prescriptions, written labels, trays, pillboxes, devices, and shared-device arrangements.

\subsubsection{Phase 2: Background and Position in the Care Network}

\begin{enumerate}
  \item \textbf{Can you tell me a little about yourself and who lives in this household?}

  \textit{Probe: Who in the household takes medicines every day, how many medicines do they take, and at what times?}
  \item \textbf{Which parts of medication management do you usually handle?}

  \textit{Probe: Walk me through the process from obtaining the medicine to the point when it is taken. Which parts are usually yours?}
  \item \textbf{How did these responsibilities come to be yours?}

  \textit{Probe: Was this discussed explicitly, or did the arrangement develop over time?}
  \item \textbf{Who else in the family takes part in this work?}

  \textit{Probe: What does each person do? Does that division remain the same or change across situations?}
\end{enumerate}

\subsubsection{Phase 3: Situated Walkthrough and Handover}

\textbf{Observational focus:} Request a physical walkthrough of a usual dose. Observe preparation, medicine identification, label reading, organization, and any improvised practices.

\begin{enumerate}[resume]
  \item \textbf{Could you show me how a typical dose is prepared and taken in this household?}

  \textit{Probe: How do you know which medicine is due and when? Show me where that information comes from.}
  \item \textbf{Can the older adult find and identify their medicines without your help?}

  \textit{Probe: Which medicines, and how do they recognize them?}
  \item \textbf{Have you used a pillbox, chart, written list, phone alarm, or another system to organize the medicines?}

  \textit{Probe: What happened when you used it? If you stopped, why?}
  \item \textbf{If you had to be away, how would this medication work be handed over to someone else?}

  \textit{Probe: Tell me about the last time this happened. What was difficult to communicate or coordinate?}
\end{enumerate}

\subsubsection{Phase 4: Disruptions, Absence, and Missed Doses}

\textbf{Observational focus:} Ground questions in recent episodes and identify where the caregiver was and what other household circumstances were present.

\begin{enumerate}[resume]
  \item \textbf{Has a dose ever been missed or substantially delayed?}

  \textit{Probe: Tell me about the most recent time. What was happening that day, and where were you?}
  \item \textbf{How did you find out that the dose had been missed, and what did you do?}

  \textit{Probe: Did you contact the older adult or someone else? What happened afterward?}
  \item \textbf{Has a missed dose ever led to a situation that concerned you?}

  \textit{Probe: Could you walk me through what happened?}
  \item \textbf{When is medication management most difficult for you?}

  \textit{Probe: What happens during work hours, travel, disrupted sleep, guests, social occasions, or religious observance?}
  \item \textbf{Has the older adult ever taken the wrong medicine or taken one twice?}

  \textit{Probe: How was this discovered, and what happened next?}
\end{enumerate}

\subsubsection{Phase 5: Prescription Interpretation and Medicine Knowledge}

\begin{enumerate}[resume]
  \item \textbf{Who usually reads and interprets prescriptions in this household?}

  \textit{Probe: Which parts can the older adult manage independently, and which parts require help?}
  \item \textbf{What do you do when something on a prescription is unclear?}

  \textit{Probe: Tell me about the last time. Did you contact a doctor, pharmacist, or another family member?}
  \item \textbf{When a new medicine is added, how do other people in the household learn about it?}

  \textit{Probe: What information do you usually explain: the medicine name, timing, purpose, or something else?}
  \item \textbf{Has a medicine ever caused a reaction or side effect that concerned you?}

  \textit{Probe: How was it identified, and whom did you consult?}
\end{enumerate}

\subsubsection{Phase 6: Technology, Proxy Use, and Device Sharing}

\textbf{Observational focus:} Where appropriate, ask the caregiver to demonstrate relevant devices or tools. Attend to who owns, holds, configures, and operates each device.

\begin{enumerate}[resume]
  \item \textbf{What phones or other relevant devices are used in this household, and whose are they?}

  \textit{Probe: Are any of the devices shared?}
  \item \textbf{What does the older adult usually use their phone for, and how did they learn to use it?}
  \item \textbf{Who usually sets up applications, alarms, or settings on household members' phones?}

  \textit{Probe: Tell me about the last thing you configured for someone else.}
  \item \textbf{Have you ever set something up for the older adult that they later stopped using, turned off, or ignored?}

  \textit{Probe: What happened?}
  \item \textbf{Are there difficulties with reading medicine information, small text, or operating the device?}

  \textit{Probe: Could you show me an example using the medicine or device you normally use?}
\end{enumerate}

\subsubsection{Phase 7: Checking, Visibility, and Decision-Making}

\begin{enumerate}[resume]
  \item \textbf{On an ordinary day, how do you know whether the older adult has taken their medicines?}

  \textit{Probe: Do you ask, observe, check the medicine, count remaining doses, or rely on another method? Tell me about yesterday.}
  \item \textbf{Does the older adult usually know when you are checking?}

  \textit{Probe: How does that checking normally happen?}
  \item \textbf{How does the older adult respond when you ask whether a medicine has been taken?}

  \textit{Probe: Can you tell me about a recent occasion?}
  \item \textbf{What is it like for you when you do not know whether a dose has been taken?}

  \textit{Probe: What do you usually do in that situation?}
  \item \textbf{When something about the medicines needs to change, how is the decision usually made?}

  \textit{Probe: Tell me about the most recent change. Who raised it, who discussed it, and who acted on it?}
  \item \textbf{Has the older adult ever changed something about the medication routine without telling you, or have you ever made a decision they did not initially agree with?}

  \textit{Probe: What happened afterward?}
  \item \textbf{Do relatives, friends, neighbours, or others outside the immediate care arrangement give medication advice?}

  \textit{Probe: Tell me about a specific occasion and what happened with that advice.}
\end{enumerate}

\subsubsection{Phase 8: Closing and Design Reflection}

\begin{enumerate}[resume]
  \item \textbf{If something could take one part of this medication work off you, which part would you most want to give up?}

  \textit{Probe: Which part would you prefer to continue handling yourself?}
  \item \textbf{If a medication reminder system existed, who should it communicate with, and in what order?}

  \textit{Probe: Should it communicate differently with the older adult and with you?}
  \item \textbf{If a scheduled dose was not confirmed, what should happen next?}

  \textit{Probe: How long should the system wait? Should another person be informed?}
  \item \textbf{Who should be able to change the system's settings?}

  \textit{Probe: If the older adult changed something, would you want to know? If you changed something, should the older adult know?}
  \item \textbf{Is there anything such a system should never do in this household?}
  \item \textbf{Is there anything about looking after the older adult's medicines that we have not discussed but that you think is important?}

  \textit{Probe: What would you want people designing medication-support technology to understand about this work?}
\end{enumerate}